\section{Theoretical Grounding}

\subsection{Self-Determination Theory}
Self-Determination Theory (SDT), developed by Ryan and Deci, proposes that intrinsic motivation is driven by three basic psychological needs: autonomy, competence, and relatedness~\cite{ryan2000self}. In educational contexts, this implies that learning systems should offer meaningful choices, provide challenges matched to the learner's ability, and foster a sense of connection or support rather than isolation.

\subsection{Activity Theory}
Activity Theory, rooted in the work of Vygotsky, provides a framework for analyzing human activity as a system involving a subject, an object, and mediating instruments, situated within a community governed by rules and division of labor~\cite{vygotsky1978mind}. It is particularly useful for identifying contradictions within a system, such as when the demands on a role exceed what one person can provide, and for reasoning about how new tools or artifacts might address those contradictions.

\subsection{Active Learning}
Active Learning refers to instructional approaches that require learners to engage in meaningful cognitive activities, such as problem-solving and reflection, rather than passively receiving information~\cite{nirenburg2017cognitive, ricci2022go}. In the context of AI-assisted education, this principle suggests that systems should promote reasoning and effort rather than simply delivering answers.

\subsection{MiniGuide}
The MiniGuide strategy, developed at Umeå University, is a pedagogical approach in which a tutor supports a learner's problem-solving by posing guiding questions rather than providing direct answers~\cite{lithner2025miniguide}. It operationalizes the concept of scaffolding within the Zone of Proximal Development, aiming to help learners develop their own reasoning rather than creating dependence on external solutions.